\section{الفصل~\ref{sec:serocomputer}. عصبونات \nt{SERO}}

خصائص عصبونات \nt{SERO} نفسها (الإيقاع، والتثبيط الذاتي، والإشارة بحسب المستقبل، والأنظمة الفرعية) مُقاسة مباشرة؛ وحسابات الفصل، أي المنظِّم والمرشّح والمنطق، تنتج من معادلاته؛ أما $\tau_c$ ومعامل انتشار السيروتونين وعدد مواقع التحرير فتقديرات؛ ودور الحلقة منظِّمًا للمستوى في الدماغ فرضية (ففي نواة الرفاء يكون التثبيط الذاتي نقطيًا ولا يعطي إلا توقفات قصيرة)، ودور \nt{SERO} أفقًا فرضية تدعمها تجارب سلوكية.

{\footnotesize
\setlength{\tabcolsep}{3pt}\renewcommand{\arraystretch}{1.15}
\begin{longtable}{>{\raggedright}p{0.5cm}>{\raggedright}p{4.0cm}>{\raggedright}p{5.6cm}>{\raggedright}p{2.3cm}>{\raggedright\arraybackslash}p{1.7cm}}
\toprule
م & القضية & التجربة وما قيس & المصدر & الحالة \\
\midrule\endfirsthead
\toprule
م & القضية & التجربة وما قيس & المصدر & الحالة \\
\midrule\endhead
1 & إشارة أثر السيروتونين يختارها مستقبِل المتلقي (القضية~\ref{prop:receptor}) & تُقسم مستقبلات السيروتونين إلى عائلات بحسب علم الأدوية والآلية: \foreignlanguage{english}{5-HT$_{1A}$} يفتح عبر بروتين \foreignlanguage{english}{G} قنوات البوتاسيوم ويثبّط الخلية، و\foreignlanguage{english}{5-HT$_{2A}$} و\foreignlanguage{english}{5-HT$_3$} المؤيني التوجه يثيران. الناقل الواحد يعطي استجابات متعاكسة في خلايا مختلفة. & \cite{BarnesSharp1999} & مُقاس \\
2 & عصبون \nt{SERO} في اليقظة يطلق بانتظام بضع نبضات في الثانية & تسجيل عصبونات نواة الرفاء الظهرية عند قطط تتحرك بحرية: تفريغ إيقاعي بطيء $2.8\pm0.2$ نبضة في الثانية في اليقظة الهادئة؛ في النوم البطيء ينخفض التردد بـ$34\text{--}68\,\%$، وفي النوم السريع بـ$98\,\%$. وتحت التخدير عند الجرذان $1.7\pm0.5$ نبضة في الثانية، بانتظام شديد. & \cite{TrulsonJacobs1979, GagnonParent2014, JacobsAzmitia1992} & مُقاس \\
3 & عصبون \nt{SERO} ناظمة إيقاع: لا يحتاج إلى دفعة لكل نبضة، لكنه يحتاج إلى دفق خلفي ثابت (في الشريحة نورأدرينالين عبر مستقبلات $\alpha_1$، وفي الكائن الحي من \nt{NORA}) & في شريحة الدماغ تفرّغ الخلايا ببطء وانتظام، وبعد كل نبضة فرط استقطاب لاحق مدته $200\text{--}800\ms$. الخلايا العاملة تلقائيًا في الشريحة أقل منها في الكائن الحي: الإيقاع المنتظم يحتاج إلى دخل نورأدرينالين مستمر عبر مستقبلات $\alpha_1$، وحجبها يخمده. & \cite{VandermaelenAghajanian1983} & مُقاس \\
4 & تكاد كل المستقبلات تكون أيضية التوجه؛ والاستجابة بطيئة: ترتفع وتهبط في غضون ثانية تقريبًا & كل عائلات المستقبلات ما عدا \foreignlanguage{english}{5-HT$_3$} مرتبطة ببروتينات \foreignlanguage{english}{G}. في الشريحة يعطي التحرير المستثار للسيروتونين عبر \foreignlanguage{english}{5-HT$_{1A}$} تيارًا مثبّطًا يرتفع ويهبط في غضون ثانية. & \cite{BarnesSharp1999, CourtneyFord2016} & مُقاس \\
5 & في مناطق الهدف يخرج السيروتونين إلى الحجم، لا إلى الشق وحده؛ وفي نواة الرفاء نفسها يجري تثبيط الجيران من نقطة إلى نقطة & قياس الفولتية الدوري السريع في الشرائح: التحرير لكل نبضة واحد في منطقة فيها مشابك وفي نواة الرفاء التي تكاد تخلو منها؛ ولا يعتمد على حجب الاسترداد والمستقبلات؛ وعدد الجزيئات لكل نهاية أقل بكثير من عدد بروتينات النقل والمستقبلات، والتركيز خارج الخلايا يطابق ألفة المستقبل الرئيسي. في نواة الرفاء يجري التثبيط عبر \foreignlanguage{english}{5-HT$_{1A}$} من نقطة إلى نقطة: بروتين النقل لا يدع السيروتونين يتراكم خارج الشق. & \cite{Bunin1998, CourtneyFord2016} & مُقاس \\
6 & التركيز وفق~\eqref{eq:seroneuron} يتناقص بسبب الضخ العائد؛ و$\tau_c$ من رتبة الثانية تقدير & قياس الفولتية في الدماغ الحي: السيروتونين خارج الخلايا يحدّه قبل كل شيء الاسترداد والتفكيك، لا التخليق. لا قياس مباشرًا لـ$\tau_c$ في الأعمال المستشهد بها؛ ورتبة المقدار تقدير الفصل. & \cite{Hashemi2012serotonin} & مُقاس؛ $\tau_c$ تقدير \\
7 & \foreignlanguage{english}{NOT} و\foreignlanguage{english}{NOR} بناظمة إيقاع واحدة؛ اكتمال منطق \nt{SERO} (القضية~\ref{prop:serocomplete}، الشكل~\ref{fig:serogates}) & ينتج من~\eqref{eq:seroneuron}: ناظمة الإيقاع القابلة للتثبيط بوابة \foreignlanguage{english}{NOR}، و\foreignlanguage{english}{NOR} كاملة وظيفيًا. لا توجد تجربة بُني فيها منطق من عصبونات \nt{SERO}؛ وشرط الأحجام المنفصلة لا يتحقق في الدماغ. & اشتقاق في الفصل & نظرية \\
8 & السيروتونين يثبّط عصبون \nt{SERO} نفسه عبر مستقبلات عليه & عند جرذان مخدّرة تثبّط ناهضات \foreignlanguage{english}{5-HT$_{1A}$}، وريديًا وبالرحلان الأيوني، تفريغ عصبونات الرفاء الظهرية بعلاقة الجرعة والاستجابة نفسها التي للسيروتونين؛ أما ناهضات \foreignlanguage{english}{5-HT$_{1B}$} فتكاد لا تؤثر. وفي الشريحة يعطي السيروتونين الداخلي المنشأ من الخلايا المجاورة عبر \foreignlanguage{english}{5-HT$_{1A}$} تيارًا وتوقفًا قصيرًا في التفريغ دون أن يغيّر التردد المتوسط. & \cite{Sprouse1987, CourtneyFord2016} & مُقاس \\
9 & في النموذج، الحلقة عبر الحجم المشترك منظِّم مستوى: التشتت والثابت الزمني يصغران بـ$1+G$ مرة~\eqref{eq:feedback}؛ وأن الحلقة تعمل هكذا في الدماغ فرضية & الصيغة الكلاسيكية للمضخّم بتغذية راجعة سالبة، وقد اشتُقت في الفصل من جديد. كسب الحلقة $G$ لنظام \nt{SERO} غير مقيس. والمقيس: في الشريحة يجري التثبيط الذاتي من نقطة إلى نقطة، ويعطي توقفًا قصيرًا ولا يغيّر التردد المتوسط. & \cite{Black1934feedback, CourtneyFord2016} & نظرية؛ وفي الدماغ فرضية \\
10 & التركيز متوسط متحرك للتردد، أي مرشّح تمرير منخفض~\eqref{eq:lowpass}، الشكل~\ref{fig:lowpass} & حل المعادلة الخطية~\eqref{eq:seroneuron}؛ والشكل حساب بهذه الصيغة عند $\tau_c=1\,\text{s}$. لا يوجد نموذج مستقل في \texttt{models/}. & اشتقاق في الفصل & نظرية \\
11 & تعرّج الشقوق يبطئ انتشار الجزيئة الصغيرة نحو $2.6$ مرة & طريقة المصدر النقطي: رحلان أيوني لرباعي ميثيل الأمونيوم وتسجيل تركيزه بقطب انتقائي للأيون في الشرائح وفي الدماغ الحي. التعرّج $\lambda\approx1.6$، أي أن $D$ الفعّال أصغر من الحر بـ$\lambda^2\approx2.6$ مرة؛ وحصة الحجم بين الخلايا نحو $0.2$. معامل انتشار السيروتونين نفسه في النسيج لم يُقَس في هذه الأعمال؛ وهو في الفصل تقدير. & \cite{Sykova2008} & مُقاس \\
12 & طول ”السلك“ المستقل $\ell\sim\sqrt{D\tau}\approx20\um$~\eqref{eq:diffusionlength}؛ وعدد الأسلاك محدود بعدد هذه المناطق (القضية~\ref{prop:volumewires}) & قانون الانتشار مع التعويض بتقديرات $D$ و$\tau$ (السطران~6 و~11)؛ والقضية~\ref{prop:volumewires} نتيجة. لا توجد تجربة مباشرة تعدّ القنوات المستقلة للنقل الحجمي. & اشتقاق في الفصل & نظرية \\
13 & خرج عصبون \nt{SERO} الواحد منتشر على مساحات هائلة من الدماغ؛ ومواقع التحرير تُقدَّر بـ$\sim10^5$ & إعادة بناء كاملة لعصبونات مفردة من الرفاء الظهرية عند الجرذ: كل المحاوير الصاعدة كثيرة التفرع، والطول الكلي لمحوار الخلية الواحدة يبلغ $18.7\,\text{cm}$، وفروع الخلية الواحدة تبلغ المخطط والقشرة الجبهية الأمامية واللوزة. عدد مواقع التحرير لكل خلية لم يُعَدّ مباشرة. & \cite{GagnonParent2014} & مُقاس؛ $10^5$ تقدير \\
14 & تنقسم عصبونات \nt{SERO} إلى بضعة أنظمة فرعية بمخارج واستجابات مختلفة & وسم جيني فيروسي عند الفئران: العصبونات الذاهبة إلى اللوزة وتلك الذاهبة إلى القشرة الجبهية تكاد لا تتقاطع في تفرعها، وتتلقى مداخل مختلفة، وتستجيب لمنبّه مزعج استجابتين متعاكستين؛ وتشغيل كل مجموعة وإيقافها يغيّران السلوك تغييرًا مختلفًا. & \cite{Ren2018} & مُقاس \\
15 & شرط الصبر $D<\tau_\gamma\ln(R/r_0)$~\eqref{eq:patience} مع الخصم الأسّي؛ ومع الخصم المقيس، الزائدي، $D<\tau_\gamma(R/r_0-1)$ & ينتج من الخصم $e^{-D/\tau_\gamma}$ أو $1/(1+D/\tau_\gamma)$؛ والاشتقاق في الفصل. قرود، اختيار على تأخيرات من رتبة الثواني: الخصم أقرب إلى المنحنى الزائدي منه إلى الأسّي؛ واستجابة عصبونات \nt{DOPA} للإشارة المنبئة بمكافأة مؤجلة تنخفض مع التأخير زائديًا أيضًا. & اشتقاق في الفصل؛ \cite{KobayashiSchultz2008} & نظرية \\
16 & \nt{SERO} يضبط الأفق $\tau_\gamma$ (القسم~\ref{sec:horizon}) & التشغيل البصري الجيني لعصبونات \nt{SERO} في الرفاء الظهرية عند الفئران يطيل انتظار المكافأة المؤجلة ويزيد المثابرة في البحث عن مكافأة صارت نادرة. وعند البشر يجعل خفضُ السيروتونين (نظام غذائي بلا تريبتوفان) خصمَ المكافأة المؤجلة أشد انحدارًا. كيف يتحول التركيز إلى $\tau_\gamma$ غير معروف. & \cite{Doya2002, Miyazaki2014, Lottem2018, Schweighofer2008} & فرضية \\
\bottomrule
\end{longtable}}
